%% secoes/introducao/01-baixo-peso.tex — incluído por main.tex

%\subsection{Baixo peso ao nascer}\label{subsec:intro-bpn}

Define-se baixo peso ao nascer como o peso inferior a 2.500 gramas ao nascimento, independentemente da idade gestacional \cite{blencowe2019}.
O baixo peso ao nascer constitui um dos mais importantes problemas de saúde pública perinatal, por refletir as condições em que a gestação transcorreu e por anteceder desfechos adversos que se estendem muito além do nascimento \cite{falcao2020,kramer1987}. Suas causas são multifatoriais e fortemente ligadas à vulnerabilidade social, o que o torna, em grande medida, evitável por meio de cuidados adequados ao longo da gestação \cite{kramer1987,fonseca2014}. Ainda assim, apesar dos avanços em saúde materno-infantil, sua ocorrência persiste elevada, no mundo e no Brasil, sem redução expressiva nas últimas décadas, configurando uma iniquidade que recai de forma desproporcional sobre as populações socioeconomicamente mais vulneráveis \cite{falcao2020,blencowe2019}.

%Define-se baixo peso ao nascer como o peso inferior a 2.500 gramas ao nascimento, independentemente da idade gestacional \cite{blencowe2019}. Suas repercussões não se restringem ao período neonatal: além do maior risco de morbidade e mortalidade nos primeiros anos de vida, indivíduos que nascem com baixo peso apresentam, na vida adulta, maior predisposição a doenças crônicas não transmissíveis, como diabetes, hipertensão e doenças cardiovasculares, associação sustentada pela hipótese das origens desenvolvimentistas da saúde e da doença \cite{barker2007}. Trata-se, portanto, de um agravo cujas consequências acompanham o indivíduo ao longo de todo o curso de vida.

%Do ponto de vista biológico, o baixo peso ao nascer decorre de prematuridade, de restrição de crescimento intrauterino ou da combinação de ambos \cite{kramer1987}. Esses mecanismos, por sua vez, são condicionados por um amplo conjunto de determinantes: características maternas, como idade e história obstétrica; a assistência pré-natal, mensurada pelo número de consultas; e fatores socioeconômicos, como escolaridade materna, renda, condições domiciliares e região de residência \cite{falcao2020}. Essa cadeia de determinação evidencia por que o desfecho se concentra justamente nos contextos de maior privação social \cite{falcao2020}.

%Por não existir medida curativa após o nascimento, o enfrentamento do baixo peso ao nascer é essencialmente preventivo. A assistência pré-natal adequada, com número suficiente de consultas e início precoce, figura entre as estratégias mais efetivas, ao permitir o acompanhamento nutricional da gestante, a identificação oportuna de complicações e a orientação quanto a comportamentos de risco \cite{falcao2020,fonseca2014}. A esse cuidado somam-se políticas de proteção social, como os programas de transferência de renda, que atuam sobre as condições socioeconômicas subjacentes ao problema \cite{falcao2020,rasella2013}.

%Quanto à magnitude, estimou-se que, em 2015, 14,6\% dos nascimentos no mundo, cerca de 20,5 milhões, ocorreram com baixo peso, dos quais 91\% em países de baixa e média renda \cite{blencowe2019}. No Brasil, a prevalência foi estimada em aproximadamente 8,5\% em 2017, patamar próximo ao observado na América Latina e no Caribe e sem tendência de queda relevante nas últimas décadas \cite{blencowe2019,ms2017}. 
%Chama atenção, no cenário nacional, o paradoxo epidemiológico do baixo peso ao nascer, em que as maiores proporções recaem sobre as regiões Sul e Sudeste, economicamente mais desenvolvidas, evidência de que a determinação do desfecho não se reduz à dimensão econômica \cite{silva2010}.
